\section{Methodology}\label{sec:methodology}

This section presents the inverse Markov-chain framework and the three
extensions developed in this work. \cref{sec:morimura-formulation}
restates the framework in our notation. \cref{sec:features}--\ref{sec:od-rebal}
present the extensions in turn.

\subsection{Parametric chain and observation model}
\label{sec:morimura-formulation}

The state space $\mathcal{X}$ is the set of directed road edges in the
largest strongly-connected component of the network. For each state
$x \in \mathcal{X}$, the next-state distribution is the softmax of
\cref{eq:softmax} over the legal successors $\mathrm{adj}(x)$, with the
score
\begin{equation}
s(x, x') \;=\; \omega^{\mathrm{loc}}_{(x,x')}
\;+\;
\underbrace{\phi_T(x')^\top \omega^{\mathrm{glo}_1}}_{\text{destination-feature term}}
\;+\;
\underbrace{\psi(x, x')^\top \omega^{\mathrm{glo}_2}}_{\text{edge-feature term}}.
\label{eq:score}
\end{equation}
Without the global terms ($d_T = d_\psi = 0$) the model has one
parameter per directed edge; with global terms it adds $d_T + d_\psi$
\emph{shared} parameters that apply across all edges of compatible type.
The initial distribution $\pI$ is similarly parametrised by per-state
weights $\nu \in \mathbb{R}^n$ via $\pI(x) \propto \exp(\nu_x)$.

We define the absorbing-corrected stationary mixture
\begin{equation}
P \;=\; \beta\,\Pchain \;+\; (1-\beta)\,\bm{1}\pI^{\!\top},
\qquad
\pi^{\!\top} P = \pi^{\!\top}, \quad \bm{1}^{\!\top}\pi = 1,
\label{eq:stationary}
\end{equation}
which models the long-run occupancy of a chain that occasionally restarts
from $\pI$ with probability $1-\beta$ per step. The hitting-rate matrix
$H(\Pchain) \in [0,1]^{|\Xo|\times|\Xo|}$ has entries equal to the
$\beta$-discounted first-passage probability between observation states,
obtained as $H_{ij} = [(I - \beta\Pchain^{(\setminus j)})^{-1}]_{ix_j}$
where $\Pchain^{(\setminus j)}$ zeroes row $j$ to encode absorption.
The Morimura observation model says we observe noisy $\fobs(x)$ for
$x \in \Xo$ and noisy $\gobs(x, x')$ for $(x, x') \in \Xo \times \Xo$,
and we wish to match the model's $\pi$ and $H$ to these observations.

\subsection{Regularised loss and analytic gradients}
\label{sec:loss-gradients}

Following \cite{morimura2013} Eqs.~7--9, the loss is the sum of a
stationary-match term, a hitting-rate-match term, and a Tikhonov ridge:
\begin{align}
L(\theta) &= \gamma\,L_d(\theta) + (1-\gamma)\,L_h(\theta) + \lambda\,R(\theta), \\
L_d(\theta) &= \tfrac{1}{2}\sum_{i, j \in \Xo}\bigl(\log\tfrac{\pi_i}{\pi_j} - \log\tfrac{f_i}{f_j}\bigr)^2, \\
L_h(\theta) &= \tfrac{1}{2}\sum_{i, j \in \Xo}\bigl(\log H_{ij} - \log g_{ij}\bigr)^2, \\
R(\theta) &= \tfrac{1}{2}\|\theta\|^2.
\end{align}
Both data terms are log-ratio losses, which is what makes them
optimisable without normalising constants on the observations
\cite{morimura2013}: any multiplicative scale on $\fobs$ or $\gobs$
cancels in the difference. The gradients are derived using the
identity for the derivative of $\log\pi$ under a softmax-parametric
chain (paper Eq.~12) and the derivative of $\log H_{ij}$ under an
absorbing-chain solve (Eq.~15). The paper's local-only gradients
$\partial L / \partial \nu$ and $\partial L / \partial \omega^{\mathrm{loc}}$
are reproduced unchanged here; we derive the omega-global
gradients (paper Eqs.~17--18 \emph{state} the parametric extension but
do not derive the gradients) by writing the chain-rule contribution of
$\omega^{\mathrm{glo}_1}$ and $\omega^{\mathrm{glo}_2}$ through the
score \eqref{eq:score} and chaining through $\nabla_\theta \log \pi$ and
$\nabla_\theta \log H$.

We verify the analytic gradients by centred finite differences
(\code{fd\_check\_features.py}) at random parameter draws across
all four blocks ($\nu$, $\omega^{\mathrm{loc}}$, $\omega^{\mathrm{glo}_1}$,
$\omega^{\mathrm{glo}_2}$). The check uses a noise-aware tolerance
($\code{rtol}=10^{-3}$ relative \emph{or} $\code{atol}=5 \times 10^{-2}$
absolute, an \code{np.allclose}-style criterion) to handle the
machine-precision FD round-off floor that dominates near-zero gradient
components. The bigger-bbox SUMO graph passes at machine precision
($\sim 10^{-8}$ relative error) on the clearly non-zero components.

\subsection{Extension 1: real edge features from the network topology}
\label{sec:features}

The Morimura framework treats $\phi_T$ and $\psi$ as user-supplied global
features. In simulated $\S 6.1$ both are drawn iid
$\mathcal{N}(0,1)$ in their setup; in their real-world Nairobi case
they use network attributes (road class, segment length, number of
lanes) extracted from the underlying OSM data.

We follow the latter convention. From the SUMO network's
\code{.net.xml}, we extract the following per-edge attributes and assemble
them into $\phi_T$ and $\psi$:
\begin{itemize}
\item \textbf{Destination-state features $\phi_T(x') \in \mathbb{R}^{8}$.}
For each destination edge $x'$: a $6$-dimensional one-hot road-class
indicator (motorway / primary / tertiary / residential / service / other),
plus the standardised log of the speed limit and the standardised log of
the lane count.
\item \textbf{Directed-edge features $\psi(x, x') \in \mathbb{R}^{2}$.}
For each directed edge $(x, x')$: the cosine of the turn angle at the
intersection (computed from the polyline geometries of $x$ and $x'$),
and the standardised log ratio of the destination speed limit to the
source speed limit. Both components are standardised to zero mean and
unit variance across the edge set.
\end{itemize}
With these definitions $d_T + d_\psi = 10$, so feature inclusion adds
just $10$ \emph{shared} parameters across the entire chain. The
expectation is that this acts as a strong regularising prior: edges of
the same road class share a common preference component, even if their
local $\omega^{\mathrm{loc}}_{(x, x')}$ weights are estimated noisily.

\subsection{Extension 2: empirical 2-step Markov diagnostic}
\label{sec:2step-diag}

The $1$-step softmax chain in \eqref{eq:softmax} embeds the standard
Markov assumption that the next edge depends only on the current edge.
A natural concern is that this assumption is empirically violated:
real drivers' next-edge choice may depend on \emph{both} the current
edge and the previous edge (e.g.\ avoiding a U-turn, preserving heading,
continuing on a high-priority arterial). If so, no fitting effort under
the $1$-step model class can close the gap to the full-observation
empirical $1$-step reference.

We test this by constructing an empirical $2$-step detector with no
fitting. From the training trajectories of each regime we count
$(x_{t-1}, x_t, x_{t+1})$ triples and form the conditional probability
$\widehat{p}_{2}(x_{t+1} \mid x_t, x_{t-1})$ as the Laplace-smoothed
empirical ratio. The detector then scores a test trajectory $\tau$ as
\[
\Lam_{2}(\tau) \;=\; \sum_{t}\bigl[\log\widehat{p}_{2}^{(B)}(\tau_{t+1}\mid\tau_t,\tau_{t-1}) - \log\widehat{p}_{2}^{(A)}(\tau_{t+1}\mid\tau_t,\tau_{t-1})\bigr],
\]
falling back to the $1$-step empirical conditional on the first step or
any unseen $(x_{t-1}, x_t)$ context. The AUC of $\Lam_{2}$ vs.\ the
$1$-step empirical $\Lam_{1}$ tells us whether higher-order memory
exists in the data. Crucially, the $2$-step detector is built from the
same trajectories and uses the same training/test split, so any
AUC difference is attributable to memory order alone.

\subsection{Extension 3: OD-rebalancing for real-world regime contrast}
\label{sec:od-rebal}

In controlled simulation, the intrinsic and sat-navigation regimes
share the same OD pairs by construction (\cref{sec:sumo}): the only
difference is the routing rule. In real-world AVI data this control is
not available: rush-hour drivers and off-peak drivers are
\emph{different populations} with different origin/destination
distributions, and any chain difference between them reflects both
their different OD mix and their different routing-given-OD choices. The
former is uninteresting (commuters vs.\ discretionary trips have
different origins by definition); the latter is what our framework is
designed to detect.

To control for OD mix we introduce an OD-rebalancing step in the
real-world loader. The procedure is:
\begin{enumerate}
\item Cluster the network's junctions into $K$ spatial zones via
$k$-means on their $(x, y)$ coordinates (standardised). Each edge is
labelled by its from-zone and its to-zone.
\item For each trip, compute its origin zone (from-zone of the first
edge in its trajectory) and its destination zone (to-zone of the last
edge).
\item Bin all regime-A trips and all regime-B trips by $(\text{zone}_o, \text{zone}_d)$.
\item For each shared OD cell, take $\min(n^A_{od}, n^B_{od})$ trips
from each regime (sampled without replacement). The matched pools
have identical marginal $(\text{zone}_o, \text{zone}_d)$ distribution.
\end{enumerate}
Any chain contrast measured on the rebalanced pools is then
attributable to the within-OD routing-choice difference, not to an
OD-mix shift. The granularity $K$ is a tuning parameter: smaller $K$
gives tighter sample sizes per cell but coarser matching; larger $K$
gives finer matching but smaller per-cell counts. We use $K=8$ as a
default (yields $\le 64$ OD cells, of which $\sim 55$ are typically
shared with sufficient counts).

\subsection{The detector and per-row diagnostic}
\label{sec:detector}

Once two chains $\Pchain^{(A)}, \Pchain^{(B)}$ are fitted, the
per-trajectory log-likelihood-ratio detector scores
\[
\Lam(\tau) \;=\; \sum_{t=0}^{T-1}\bigl[\log\pT^{(B)}(\tau_{t+1}\mid\tau_t) - \log\pT^{(A)}(\tau_{t+1}\mid\tau_t)\bigr],
\]
and the AUC of $\Lam$ against the labelled test trajectories
($N$ per regime, $T$-edge windows) summarises detection performance.
Following the iter-18 finding that AUC alone can reflect aggregate
chain-drift rather than per-row chain accuracy, we additionally report
a \emph{contrast-correlation} diagnostic at branching states (states
with out-degree $> 1$): for each such state $x$, we compute
\[
\begin{aligned}
\Delta_{\mathrm{fit}}(x, \cdot) &= p^{(B)}_{T,\mathrm{fit}}(x, \cdot) - p^{(A)}_{T,\mathrm{fit}}(x, \cdot), \\
\Delta_{\mathrm{emp}}(x, \cdot) &= p^{(B)}_{T,\mathrm{emp}}(x, \cdot) - p^{(A)}_{T,\mathrm{emp}}(x, \cdot),
\end{aligned}
\]
where the empirical chain $P_{T,\mathrm{emp}}$ is the
full-observation Laplace-smoothed transition matrix from the trajectories.
We report the stacked Pearson correlation of $\Delta_{\mathrm{fit}}$
against $\Delta_{\mathrm{emp}}$ over all branching-state legal-edge entries,
plus the per-state cosine similarity distribution (median, mean, and
fraction of states with positive cosine). A high AUC with low Pearson
indicates aggregate-bias-driven detection; a high AUC \emph{and} a
non-trivial Pearson together support the stronger ``chain recovery''
claim.

\subsection{Design rationale: each extension targets a specific failure mode}
\label{sec:design-rationale}

The extensions above are not independent features but a deliberate
sequence, each addressing a failure mode exposed by the previous one:
\begin{itemize}
\item \textbf{Real edge features} (\cref{sec:features}) address sparse
observation and parameter explosion: $10$ shared parameters let edges of
the same road class borrow strength, regularising the per-edge weights
that are otherwise estimated from few or no observed transitions.
\item \textbf{Per-row Pearson / cosine diagnostics} (\cref{sec:detector})
address the risk that AUC improves for aggregate-bias reasons rather than
correct chain recovery, by measuring the fitted-vs-empirical contrast
directly at branching states.
\item \textbf{The empirical $2$-step diagnostic} (\cref{sec:2step-diag})
tests whether higher-order memory is an obvious missing lever before any
parametric $2$-step model is built.
\item \textbf{OD-rebalancing} (\cref{sec:od-rebal}) addresses the
real-world population-mix confound that has no analogue in the controlled
SUMO comparison.
\item \textbf{Multistart / basin checks} (\cref{sec:phase4-multistart},
\cref{sec:phase5-multistart}) address the possibility that a result is an
L-BFGS optimiser artefact rather than a property of the data.
\end{itemize}
Where AUCs lie close to chance, the honest interpretation also requires
quantifying sampling noise; we return to this as a limitation
(\cref{sec:phase5-reading}, \cref{sec:disc-limits}) since the bootstrap
confidence intervals it calls for were not computed in the time available.
